\section*{Impact Statement}

This work introduces a framework for the rigorous behavioral profiling of Large Language Models (LLMs). As LLMs increasingly transition from passive text generators to active agents involved in decision-support, prioritization, and recommendation systems, their latent behavioral tendencies, such as the ethnocentric biases analyzed in our case study, can have tangible downstream consequences. 
For instance, a model deployed in an e-commerce or procurement setting that systematically favors products from specific regions could introduce subtle but scalable economic distortions. 
By providing a methodology to quantify these tendencies with statistical precision, we aim to empower auditors, developers, and policymakers to detect and mitigate such biases before deployment.

While our framework adopts terminology and tools from psychology (e.g., "consensus", "constructs", "ethnocentrism"), it is crucial to interpret these results through a statistical lens rather than an anthropomorphic one. The detection of ethnocentrism in an LLM describes a distributional regularity in token prediction, not a conscious belief system or intent held by the model or its developers. We caution against the misuse of such behavioral profiles to ascribe moral agency to algorithmic systems or to fuel geopolitical narratives regarding the ``national character'' of specific models.

Furthermore, while our primary goal is bias detection and mitigation, we acknowledge that a framework for measuring latent behavioral constructs could theoretically be used to optimize models for specific, potentially manipulative, personality traits. However, we believe that the benefits of transparency (i.e., moving behavioral evaluation from anecdotal to rigorous, reproducible science) far outweigh the risks. 
By making the stochastic nature of model behavior explicit, this work contributes to a more mature and realistic understanding of AI safety and alignment.
